%=====================================================================
\chapter{Hidden Markov Models}\label{ch:hmm}
%=====================================================================
\locator{5}{Part V --- Sequences, Decisions and Reinforcement}

\begin{outcomes}
By the end of this chapter you will be able to:
\begin{tight}
  \item \textbf{State} the Markov property, and \textbf{compute} with a Markov
        chain's transition matrix.
  \item \textbf{Specify} a hidden Markov model by its states, observations, initial,
        transition and emission probabilities.
  \item \textbf{Compute} the probability of an observation sequence with the forward
        algorithm, by hand.
  \item \textbf{Find} the most likely hidden-state sequence with the Viterbi
        algorithm, by hand.
  \item \textbf{Explain} how Baum--Welch learns an HMM's parameters, and
        \textbf{estimate} a probability by Monte Carlo sampling.
\end{tight}
\end{outcomes}

\begin{prereq}
\chref{math} (conditional probability), \chref{bayes} (Bayes' theorem, posteriors)
and \chref{semisup} (the EM algorithm, which trains HMMs). The HEC outline lists HMMs
under reinforcement learning; they belong here, at the start of Part V, because the
Markov machinery they introduce is what Markov decision processes (\chref{mdp}) build
on.
\end{prereq}

\begin{keyterms}
sequence $\bullet$ Markov property $\bullet$ Markov chain $\bullet$ transition
matrix $\bullet$ stationary distribution $\bullet$ hidden state $\bullet$ observation
(emission) $\bullet$ initial distribution $\pi$ $\bullet$ emission matrix $\bullet$
evaluation, decoding and learning problems $\bullet$ forward algorithm $\bullet$
Viterbi algorithm $\bullet$ backpointer $\bullet$ Baum--Welch $\bullet$ Monte Carlo
estimation
\end{keyterms}

\section{When Order Matters}

Every model so far has treated examples as independent: shuffling the training
set changes nothing. Much real data is a \emph{sequence} in which the order carries
the information --- the words of a sentence, the readings of a sensor, the actions
of a user in a session, the calls a program makes to the operating system. Hidden
Markov models were the standard tool for such data for thirty years, in speech
recognition, part-of-speech tagging and gene finding, and they remain the clearest
introduction to reasoning about sequences.

\section{Markov Chains}

\begin{definitionbox}[Markov Property and Markov Chain]
A sequence of states $s_1, s_2, \dots$ has the \term{Markov property} if the next
state depends only on the current one, not on the history:
$P(s_{t+1} \mid s_t, s_{t-1}, \dots, s_1) = P(s_{t+1} \mid s_t)$. A \term{Markov
chain} is such a sequence, described by an initial distribution $\pi$ and a
\term{transition matrix} $A$ with $A_{ij} = P(s_{t+1} = j \mid s_t = i)$; each row
of $A$ sums to 1.
\end{definitionbox}

A server is either Healthy (H) or Degraded (D). From one minute to the next it stays
healthy with probability 0.8 and degrades with probability 0.2; once degraded, it
recovers with probability 0.3 and stays degraded with 0.7:
\[
A = \begin{pmatrix} 0.8 & 0.2 \\ 0.3 & 0.7 \end{pmatrix}
\]
If it is healthy now, the probability it is degraded in two minutes sums over the
two routes: H$\to$H$\to$D $= 0.8 \times 0.2 = 0.16$ and H$\to$D$\to$D $= 0.2 \times 0.7
= 0.14$, total 0.30 --- which is the entry of $A^2$ in row H, column D. Run long
enough, the chain forgets where it started and settles into a \term{stationary
distribution}: here 60\% healthy, 40\% degraded (check: $0.6 \times 0.2 = 0.4 \times
0.3$, so the flows between the states balance).

\section{Hidden Markov Models}

In practice the server's true state is not visible. What the monitoring sees is its
\emph{response time}, which depends on the state but only probabilistically: a
healthy server is usually fast, a degraded one usually slow or timing out, but either
can produce any reading. The states are \emph{hidden}; the observations are what is
recorded.

\begin{definitionbox}[Hidden Markov Model]
An HMM has hidden states $S$ and observable symbols $O$, and three sets of
probabilities:
\begin{tight}
  \item the \term{initial distribution} $\pi_i = P(s_1 = i)$;
  \item the \term{transition matrix} $A_{ij} = P(s_{t+1} = j \mid s_t = i)$;
  \item the \term{emission matrix} $B_{ik} = P(o_t = k \mid s_t = i)$.
\end{tight}
The hidden states form a Markov chain, and each observation depends only on the
hidden state at the same time.
\end{definitionbox}

\dsfig{%
\begin{tikzpicture}[font=\scriptsize,
  st/.style={circle, draw=#1, line width=1.1pt, fill=#1!12, text=#1!55!black,
             minimum size=1.3cm, font=\scriptsize\bfseries, align=center},
  ob/.style={rounded corners=3pt, draw=gray!55, fill=gray!8, minimum width=1.4cm,
             minimum height=0.6cm, font=\tiny\bfseries},
  a/.style={-{Stealth[length=2mm]}, line width=0.9pt}]
\node[st=greenhl!70!black] (H) at (0,0) {Healthy};
\node[st=accent] (D) at (4.6,0) {Degraded};
\draw[a, greenhl!60!black] (H) to[out=130, in=50, looseness=6] node[above, font=\tiny] {0.8} (H);
\draw[a, accent] (D) to[out=130, in=50, looseness=6] node[above, font=\tiny] {0.7} (D);
\draw[a, gray!60!black] (H) to[bend left=25] node[above, font=\tiny] {0.2} (D);
\draw[a, gray!60!black] (D) to[bend left=25] node[below, font=\tiny] {0.3} (H);
\node[ob] (f) at (-0.6,-2.9) {fast};
\node[ob] (s) at (2.3,-2.9) {slow};
\node[ob] (t) at (5.2,-2.9) {timeout};
\foreach \o in {f,s,t}{
  \draw[a, greenhl!60!black, dashed, line width=0.6pt] (H) -- (\o);
  \draw[a, accent, dashed, line width=0.6pt] (D) -- (\o);}
\node[anchor=north west, font=\scriptsize] at (7.0,1.3) {%
  \renewcommand{\arraystretch}{1.25}%
  \begin{tabular}{@{}lccc@{}}
  \multicolumn{4}{@{}l}{\textbf{Emissions} $B$: $P(\text{reading} \mid \text{state})$} \\[2pt]
  & fast & slow & timeout \\
  {\color{greenhl!45!black}\bfseries H} & 0.70 & 0.25 & 0.05 \\
  {\color{accent}\bfseries D} & 0.10 & 0.50 & 0.40 \\[4pt]
  \multicolumn{4}{@{}l}{\textbf{Start} $\pi$: H 0.7, D 0.3} \\[2pt]
  \multicolumn{4}{@{}l}{\textbf{Transitions} $A$: solid arrows} \\
  \end{tabular}};
\node[font=\tiny, text=gray!55!black, anchor=west] at (7.0,-2.9)
     {\textbf{hidden}: the state (top); \textbf{observed}: the reading (bottom)};
\end{tikzpicture}}%
{The hidden Markov model used throughout this chapter. The monitoring system never sees
the state, only a response time whose probabilities depend on it.}{hmm-model}

\begin{conceptbox}[Three Questions an HMM Answers]
\begin{tight}
  \item \textbf{Evaluation:} how probable is an observed sequence under this model?
        (Is this pattern of response times normal?) --- the forward algorithm.
  \item \textbf{Decoding:} which sequence of hidden states most probably produced
        it? (When did the server actually degrade?) --- the Viterbi algorithm.
  \item \textbf{Learning:} which parameters make the observed sequences most
        probable? --- the Baum--Welch algorithm, an instance of EM.
\end{tight}
\end{conceptbox}

\section{Evaluation: The Forward Algorithm}

The direct way to find $P(o_1, \dots, o_T)$ sums over every possible hidden-state
sequence --- $N^T$ of them, which for 2 states and 100 time steps is $2^{100}$. The
forward algorithm avoids the explosion by reusing partial sums.

\begin{definitionbox}[Forward Algorithm]
Let $\alpha_t(i) = P(o_1, \dots, o_t,\ s_t = i)$: the probability of the observations
so far \emph{and} being in state $i$ now.
\begin{tight}
  \item Start: $\alpha_1(i) = \pi_i\,B_i(o_1)$.
  \item Step: $\alpha_{t+1}(j) = \Big[\sum_i \alpha_t(i)\,A_{ij}\Big]\,B_j(o_{t+1})$
        --- every way of arriving in $j$, times the chance of emitting what was seen.
  \item End: $P(o_1, \dots, o_T) = \sum_i \alpha_T(i)$.
\end{tight}
The cost is $N^2 T$ instead of $N^T$.
\end{definitionbox}

\begin{worked}[How Likely Is ``Fast, Slow, Timeout''?]
Observations $o = $ (fast, slow, timeout), model of \dsref{hmm-model}.

\smallskip
\textbf{$t = 1$ (fast):} $\alpha_1(H) = 0.7 \times 0.70 = 0.49$;
$\alpha_1(D) = 0.3 \times 0.10 = 0.03$.

\textbf{$t = 2$ (slow):}
\[
\alpha_2(H) = (0.49 \times 0.8 + 0.03 \times 0.3) \times 0.25 = 0.401 \times 0.25 = 0.10025
\]
\[
\alpha_2(D) = (0.49 \times 0.2 + 0.03 \times 0.7) \times 0.50 = 0.119 \times 0.5 = 0.0595
\]
\textbf{$t = 3$ (timeout):}
\[
\alpha_3(H) = (0.10025 \times 0.8 + 0.0595 \times 0.3) \times 0.05 = 0.09805 \times 0.05 = 0.00490
\]
\[
\alpha_3(D) = (0.10025 \times 0.2 + 0.0595 \times 0.7) \times 0.40 = 0.0617 \times 0.4 = 0.02468
\]
\textbf{Answer:} $P(o) = 0.00490 + 0.02468 = \mathbf{0.0296}$.

\smallskip
The same numbers answer a second question for free: given everything seen so far,
the server is now degraded with probability $0.02468/0.0296 = \mathbf{0.83}$ ---
the \emph{filtered} state estimate, which is what a live monitoring dashboard wants.
\end{worked}

\section{Decoding: The Viterbi Algorithm}

Viterbi has the same structure as the forward algorithm with one change: replace the
sum over incoming paths by a \emph{maximum}, and remember which predecessor gave it.

\begin{definitionbox}[Viterbi Algorithm]
Let $\delta_t(j)$ be the probability of the single most likely state sequence ending
in state $j$ at time $t$ (with the observations so far).
\begin{tight}
  \item Start: $\delta_1(i) = \pi_i\,B_i(o_1)$.
  \item Step: $\delta_{t+1}(j) = \max_i\big[\delta_t(i)\,A_{ij}\big]\,B_j(o_{t+1})$,
        and store the maximising $i$ as the \term{backpointer} for $(t+1, j)$.
  \item End: take the state with the largest $\delta_T$, then follow the
        backpointers back to $t = 1$.
\end{tight}
\end{definitionbox}

\begin{worked}[When Did the Server Degrade?]
The same observations: fast, slow, timeout.

\smallskip
\textbf{$t = 1$:} $\delta_1(H) = 0.49$, $\delta_1(D) = 0.03$.

\textbf{$t = 2$ (slow):}
\begin{tight}
  \item into H: from H $0.49 \times 0.8 = 0.392$, from D $0.03 \times 0.3 = 0.009$;
        best from H. $\delta_2(H) = 0.392 \times 0.25 = 0.098$.
  \item into D: from H $0.49 \times 0.2 = 0.098$, from D $0.03 \times 0.7 = 0.021$;
        best from H. $\delta_2(D) = 0.098 \times 0.5 = 0.049$.
\end{tight}
\textbf{$t = 3$ (timeout):}
\begin{tight}
  \item into H: from H $0.098 \times 0.8 = 0.0784$, from D $0.049 \times 0.3 = 0.0147$;
        best from H. $\delta_3(H) = 0.0784 \times 0.05 = 0.00392$.
  \item into D: from H $0.098 \times 0.2 = 0.0196$, from D $0.049 \times 0.7 = 0.0343$;
        best from D. $\delta_3(D) = 0.0343 \times 0.4 = \mathbf{0.01372}$.
\end{tight}
\textbf{Backtrack.} The best final state is D (0.01372). Its backpointer is D at
$t = 2$, whose backpointer is H at $t = 1$. \textbf{Most likely states: H, D, D} ---
the server degraded between the first and second minute, when the first slow
response appeared, not when the timeout did.
\end{worked}

\dsfig{%
\begin{tikzpicture}[font=\scriptsize,
  nd/.style={circle, draw=#1, line width=1pt, fill=#1!10, minimum size=1.15cm, align=center,
             font=\tiny},
  e/.style={gray!40, line width=0.7pt},
  best/.style={-{Stealth[length=2.2mm]}, primary, line width=2pt}]
\foreach \t/\o in {0/fast,1/slow,2/timeout}{
  \node[font=\scriptsize\bfseries, text=gray!55!black] at (\t*3.6,1.5) {$t = \the\numexpr\t+1\relax$: \o};}
\node[font=\scriptsize\bfseries, text=greenhl!45!black, anchor=east] at (-1.0,0.5) {Healthy};
\node[font=\scriptsize\bfseries, text=accent, anchor=east] at (-1.0,-1.1) {Degraded};
\node[nd=greenhl!70!black] (h1) at (0,0.5) {0.49};
\node[nd=accent] (d1) at (0,-1.1) {0.03};
\node[nd=greenhl!70!black] (h2) at (3.6,0.5) {0.098};
\node[nd=accent] (d2) at (3.6,-1.1) {0.049};
\node[nd=greenhl!70!black] (h3) at (7.2,0.5) {0.00392};
\node[nd=accent, line width=2pt] (d3) at (7.2,-1.1) {0.01372};
\foreach \a/\b in {h1/h2,d1/h2,d1/d2,h2/h3,d2/h3,h2/d3}{\draw[e] (\a) -- (\b);}
\draw[best] (h1) -- (d2);
\draw[best] (d2) -- (d3);
\draw[gray!60, -{Stealth[length=1.8mm]}, line width=0.9pt] (h1) -- (h2);
\node[anchor=west, align=left, text width=4.6cm, font=\scriptsize, text=gray!55!black] at (8.6,-0.3)
     {Each node shows $\delta_t$: the probability of the best path ending there. The
      thick arrows are the backpointers of the winning path, H $\to$ D $\to$ D.};
\end{tikzpicture}}%
{The Viterbi trellis for the worked example. Every hidden-state sequence is a path through
the trellis; Viterbi keeps only the best path into each node.}{viterbi}

\section{Learning: Baum--Welch}

Where do $\pi$, $A$ and $B$ come from? If the hidden states had been recorded, they
would be counts: how often H followed H, how often D emitted a timeout. They were not
--- so the \term{Baum--Welch} algorithm applies EM (\chref{semisup}). The
\textbf{E-step} runs the forward algorithm and its mirror image, the backward
algorithm, to compute for every time step the probability of each hidden state and of
each transition, given the whole observed sequence. The \textbf{M-step} re-estimates
$A$ and $B$ as expected counts: the expected number of H$\to$D transitions divided by
the expected number of visits to H, and so on. As with every EM, the likelihood never
decreases and the answer is a local maximum. This chapter's lab generates 3{,}000
response times from the model, forgets the parameters, and from a vague start recovers
\[
A \approx \begin{pmatrix} 0.78 & 0.22 \\ 0.31 & 0.69 \end{pmatrix},
\]
close to the true values.

\section{Monte Carlo Estimation}

There is a completely different way to answer the evaluation question: \emph{simulate}.
Generate many sequences from the model and count how often the observed one appears.

\begin{definitionbox}[Monte Carlo Estimation]
To estimate a probability or an expectation, draw many independent samples from the
model and average: the fraction of samples in which an event occurs estimates its
probability, and the average of a quantity over the samples estimates its expected
value. The error shrinks like $1/\sqrt{\text{number of samples}}$ --- the same
square-root law as the confidence intervals of \chref{evaluation}.
\end{definitionbox}

For three observations the forward algorithm is exact and instant, and simulation is
pointless: with 1{,}000 samples the lab's estimate is 0.025, with 10{,}000 it is 0.032,
and with 100{,}000 it is 0.0296, finally matching. Monte Carlo earns its place when no
exact algorithm exists --- models too complex to sum over, or, as in \chref{rl}, an
environment whose probabilities are unknown and can only be experienced. There,
averaging the outcomes of many trials is the only estimate available, and it is the
foundation of Monte Carlo reinforcement learning.

\begin{alertbox}[Why HMMs Were Largely Replaced --- and Where They Remain]
An HMM's hidden state is one symbol from a small set, and each observation depends on
the current state alone. Language and speech have longer dependencies than that, and
recurrent networks and transformers (the Deep Learning course) now dominate both.
HMMs remain the right tool when the data is small, the states are meaningful and must
be interpretable --- a machine's health, a user's session phase, the regime of a
market --- and when an exact probability is needed rather than a score.
\end{alertbox}

\begin{pitfall}
\begin{tight}
  \item \textbf{Multiplying probabilities along long sequences.} They underflow to
        zero; scale at each step or work in logarithms (Viterbi is usually run in log
        space, turning products into sums).
  \item \textbf{Confusing the forward and Viterbi results.} Forward sums over all
        paths (probability of the observations); Viterbi keeps the single best path.
  \item \textbf{Taking the most likely state at each step as the best sequence.}
        The individually likeliest states may form a sequence that is itself
        unlikely, or even impossible; decode with Viterbi.
  \item \textbf{Trusting Baum--Welch from one start.} It finds a local maximum;
        restart from several initial guesses.
  \item \textbf{Assuming the Markov property holds.} If the next state depends on a
        long history, a first-order HMM is the wrong model.
\end{tight}
\end{pitfall}

%---------------------------------------------------------------------
\begin{fourlenses}{Hidden Markov Models}
\lensLA{A student's hidden state --- engaged, struggling, disengaged --- produces
observable behaviour: logins, quiz scores, forum posts. An HMM trained on past
cohorts decodes when a student moved from struggling to disengaged, which is when an
intervention would have helped most.}

\lensPM{A project's hidden health (on track, at risk, in trouble) emits weekly
signals: velocity, defect counts, missed stand-ups. The filtered probability of ``in
trouble'' is a more stable early warning than any single metric.}

\lensIOT{The worked example is realistic: equipment degrades through hidden stages
that show up only through noisy sensor readings. HMM-based condition monitoring and
remaining-useful-life estimation are standard in predictive maintenance.}

\lensSEC{The sequence of system calls made by a program was one of the first
intrusion-detection signals: an HMM of a normal program's calls gives an
attacked program's call sequence a very low forward probability. Session-level user
behaviour can be modelled the same way.}
\end{fourlenses}

%---------------------------------------------------------------------
\begin{lab}[Forward, Viterbi, Monte Carlo and Baum-Welch]
\begin{lstlisting}[language=Python]
import numpy as np

# --- the worked example: a server that is Healthy (0) or Degraded (1) ----
pi = np.array([0.7, 0.3])                        # initial state
A = np.array([[0.8, 0.2],                        # transitions, row = from
              [0.3, 0.7]])
B = np.array([[0.70, 0.25, 0.05],                # emissions: fast, slow,
              [0.10, 0.50, 0.40]])               #            timeout
obs = [0, 1, 2]                                  # fast, slow, timeout


def forward(obs, pi, A, B):
    alpha = pi * B[:, obs[0]]
    for o in obs[1:]:
        alpha = (alpha @ A) * B[:, o]
    return alpha                                 # sums to P(obs)


def viterbi(obs, pi, A, B):
    delta, back = pi * B[:, obs[0]], []
    for o in obs[1:]:
        cand = delta[:, None] * A                # cand[i, j]: from i to j
        back.append(cand.argmax(axis=0))
        delta = cand.max(axis=0) * B[:, o]
    path = [int(delta.argmax())]
    for b in reversed(back):
        path.insert(0, int(b[path[0]]))
    return path, delta.max()


alpha = forward(obs, pi, A, B)
print("P(fast, slow, timeout) =", round(alpha.sum(), 6))
print("P(state at t=3 | obs)  =", (alpha / alpha.sum()).round(3))
path, p = viterbi(obs, pi, A, B)
print("most likely states:", ["HD"[s] for s in path], "prob", round(p, 5))

# --- Monte Carlo: estimate the same probability by sampling --------------
rng = np.random.default_rng(0)


def sample(T):
    s = rng.choice(2, p=pi)
    out = []
    for t in range(T):
        out.append(rng.choice(3, p=B[s]))
        s = rng.choice(2, p=A[s])
    return out


for n in (1_000, 10_000, 100_000):
    hits = sum(sample(3) == obs for _ in range(n))
    print(f"Monte Carlo, {n:6d} samples: {hits / n:.4f}")

# --- Baum-Welch (EM): relearn A and B from one long sampled sequence -----
seq = np.array(sample(3000))
A_hat = np.array([[0.5, 0.5], [0.5, 0.5]])       # deliberately vague start
B_hat = np.array([[0.5, 0.3, 0.2], [0.2, 0.3, 0.5]])
p_hat = np.array([0.5, 0.5])
T = len(seq)
for it in range(60):
    a, c = np.zeros((T, 2)), np.zeros(T)         # scaled forward pass
    a[0] = p_hat * B_hat[:, seq[0]]
    c[0] = a[0].sum()
    a[0] /= c[0]
    for t in range(1, T):
        a[t] = (a[t - 1] @ A_hat) * B_hat[:, seq[t]]
        c[t] = a[t].sum()
        a[t] /= c[t]

    b = np.ones((T, 2))                          # scaled backward pass
    for t in range(T - 2, -1, -1):
        b[t] = A_hat @ (B_hat[:, seq[t + 1]] * b[t + 1]) / c[t + 1]

    g = a * b                                    # E-step: P(state_t | obs)
    xi = (a[:-1, :, None] * A_hat[None] *
          (B_hat[:, seq[1:]].T * b[1:])[:, None, :]) / c[1:, None, None]
    A_hat = xi.sum(0) / g[:-1].sum(0)[:, None]   # M-step: expected counts
    B_hat = np.array([g[seq == k].sum(0) for k in range(3)]).T
    B_hat /= g.sum(0)[:, None]
    p_hat = g[0]
print("relearned A:\n", A_hat.round(2))
print("relearned B:\n", B_hat.round(2))
\end{lstlisting}
The forward and Viterbi results reproduce both worked examples. Change the start of
Baum--Welch to \texttt{B\_hat = [[0.2, 0.3, 0.5], [0.5, 0.3, 0.2]]}: the states come out
the other way round. Why is that not an error?
\end{lab}

%---------------------------------------------------------------------
\begin{checkpoint}
\begin{tightnum}
  \item In the Markov chain of this chapter, a server is degraded now. What is the
        probability it is healthy two minutes later?
  \item What does $\alpha_t(i)$ represent, and how does $\delta_t(i)$ differ from it?
  \item Why does the direct sum over all hidden sequences become impractical, and
        how does the forward algorithm avoid it?
\end{tightnum}
\end{checkpoint}

%---------------------------------------------------------------------
\begin{chaptersummary}
\begin{tight}
  \item A Markov chain's next state depends only on the current one; its transition
        matrix gives multi-step probabilities and a stationary distribution.
  \item An HMM adds observations emitted from hidden states: $\pi$, $A$ and $B$
        define it completely.
  \item The forward algorithm computes the probability of an observation sequence in
        $N^2 T$ steps, and yields the filtered state probabilities.
  \item Viterbi replaces the forward sum by a maximum with backpointers, giving the
        most likely hidden-state sequence.
  \item Baum--Welch learns the parameters by EM; Monte Carlo sampling estimates any
        probability by simulation, with error falling as $1/\sqrt{n}$.
\end{tight}
\end{chaptersummary}

\begin{reviewq}
  \item \textbf{(MCQ)} In an HMM, the emission matrix gives: (a)~$P(s_{t+1} \mid s_t)$
        (b)~$P(o_t \mid s_t)$ (c)~$P(s_1)$ (d)~$P(s_t \mid o_t)$.
  \item \textbf{(MCQ)} The Viterbi algorithm differs from the forward algorithm by
        replacing a sum with: (a)~a product (b)~a maximum (c)~an average (d)~a
        logarithm.
  \item \textbf{(MCQ)} Baum--Welch is an instance of: (a)~gradient descent
        (b)~the EM algorithm (c)~$k$-means (d)~Q-learning.
  \item \textbf{(Short)} State the three problems an HMM can solve and the algorithm
        for each.
  \item \textbf{(Short)} Why can choosing the most likely state at each time step
        separately give a worse answer than Viterbi?
  \item \textbf{(Short)} When is Monte Carlo estimation preferable to an exact
        algorithm?
  \item \textbf{(Applied)} Compute $P(\text{slow, slow})$ for the model of this
        chapter with the forward algorithm, and the filtered probability that the
        server is degraded after the second observation.
  \item \textbf{(Applied)} Run Viterbi on the observations (timeout, fast, fast) and
        interpret the decoded state sequence.
\end{reviewq}
